\section{The simulated boat and water}
\label{sec:sim}

The marine surface profile supplies the boat's rigid inertia, added mass,
damping and individual thruster specifications. The Ideal profile uses the
following defaults; the equations are developed in Section~\ref{sec:s2r}.

\paragraph{Thrust.} Four thrusters have symmetric full-effort magnitudes of
$51.5$\,N at the reference voltage, zero command delay and zero lag. Their
signed quadratic law gives $4\cdot51.5\cos45^\circ\approx146$\,N at full surge.

\paragraph{Hull drag.} On water-relative body velocity, each horizontal axis
has linear-plus-quadratic damping. Ideal coefficients $(D_1,D_2)$ are
$(10,19.4)$ for surge, $(15,41.7)$ for sway and $(10,19.3)$ for yaw, in the
corresponding SI units, symmetric in sign. Added masses are $1.5$\,kg in surge,
$15$\,kg in sway and $1.23$\,kg\,m$^2$ in yaw, with zero sway--yaw coupling.
The profile step is $0.02$\,s.

\paragraph{Mass and inertia.} 30\,kg loaded. Two hulls carry about 60\,\% of
it as slender $1.06$\,m bodies $0.33$\,m either side of the centre line, the deck,
electronics and batteries the rest as a $0.5\times0.4\times0.25$\,m box, so about
the centre of mass
\[
  I_\text{pitch} \approx 2\cdot 9\,\tfrac{1.06^2}{12} + 12\,\tfrac{0.5^2+0.25^2}{12} = 2.0,\qquad
  I_\text{yaw} \approx 2\cdot 9\left(\tfrac{1.06^2}{12} + 0.33^2\right) + 12\,\tfrac{0.5^2+0.4^2}{12} = 4.1
\]
kg\,m$^2$. These motivate the profile's pitch and yaw inertia priors.
The profile's roll inertia is $2.2$\,kg\,m$^2$. Rotor links retain their authored
masses; fixed and root links share the remainder of the 30\,kg total. Composite
inertia uses the parallel-axis theorem,
$I=\sum_i[I_i+m_i(\|\bm d_i\|^2-(\bm d_i\cdot\bm a)^2)]$ about axis $\bm a$.
The root supplies the remaining inertia after the other links' contributions.

\paragraph{Wind.} With the wind option enabled, air at velocity $\bm w$ pushes the hull's above-water area with
$\bm F = \tfrac12 \rho_a C_d A(\beta)\,|\bm w_r|\,\bm w_r$, $\bm w_r = \bm w - \bm v$,
where $A(\beta) = A_f|\cos\beta| + A_s|\sin\beta|$ is the area facing the relative
wind at $\beta$ from the bow ($A_f = 0.30$, $A_s = 0.40$\,m$^2$, $C_d = 1$). An
$8$\,m/s breeze abeam gives $15.7$\,N. Gusts and veer follow Ornstein--Uhlenbeck
processes, $\dot g = -g/T + \sigma\sqrt{2/T}\,\xi$: speed varies by $25\,\%$ over
about $8$\,s, direction by $10^\circ$ over $30$\,s. The force acts above and
slightly ahead of the centre of mass, so the boat weathervanes.

\paragraph{Sensors.} The LiDAR is a VLP-16 in its own REP-103 frame (16 rings,
$2^\circ$ apart over $\pm 15^\circ$, $0.2^\circ$ steps, 10\,Hz), publishing only
returns, and none from below the water surface. Detections are in the camera's
REP-103 frame. Both must line up with the boat description's mounts: a sensor
$0.47$\,m off moves every buoy by as much, about the clearance in a gate.
\file{crane\_ml: MarineSurfaceProfile.cs, RoboBoatPhysics.cs, MarineHydrodynamics.cs, ThrusterModel.cs, RoboBoatWind.cs, RoboBoatCourse.cs, Lidar3D.cs}
